% =====================================================================
%  Klimabonus Frankfurt — Feldexperiment
%  Studien-Briefing für die Pre-Registrierung
%
%  Zweck: Dieses Dokument fasst Design, erfasste Variablen und erste
%  Hypothesen-Skizze für die Pre-Reg-Vorbereitung durch die Co-Autoren
%  zusammen. Es ist kein finaler Pre-Analysis-Plan, sondern eine
%  strukturierte Übersicht.
%
%  Kompilation:  pdflatex study_briefing.tex
% =====================================================================

\documentclass[11pt, a4paper]{scrartcl}

\usepackage[utf8]{inputenc}
\usepackage[T1]{fontenc}
\usepackage[ngerman]{babel}
\usepackage[margin=2.4cm]{geometry}
\usepackage{lmodern}
\usepackage{microtype}
\usepackage{xcolor}
\usepackage{booktabs}
\usepackage{longtable}
\usepackage{tabularx}
\usepackage{hyperref}
\usepackage{enumitem}
\usepackage{xurl}

\definecolor{gublue}{HTML}{003865}
\definecolor{gugrey}{HTML}{5a6573}
\hypersetup{colorlinks=true, urlcolor=gublue, linkcolor=gublue, citecolor=gublue}
\setlist{noitemsep, topsep=2pt}

\setkomafont{title}{\color{gublue}\bfseries}
\setkomafont{section}{\color{gublue}\Large\bfseries}
\setkomafont{subsection}{\color{gublue}\large\bfseries}

\renewcommand{\arraystretch}{1.15}
\newcommand{\var}[1]{\texttt{\small #1}}

\title{Klimabonus Frankfurt --- Feldexperiment\\
\Large Studien-Briefing für die Pre-Registrierung}
\author{Goethe-Universität Frankfurt \(\cdot\) ZEW Mannheim \(\cdot\)
        Klimareferat Stadt Frankfurt}
\date{Stand: \today}

\begin{document}
\maketitle

% =====================================================================
\section*{Executive Summary}
% =====================================================================

Drei-armiges randomisiertes Feldexperiment zum Förderprogramm
\emph{Klimabonus} der Stadt Frankfurt. Etwa 20.000 Frankfurter Unternehmen
erhalten einen \textbf{treatment-neutralen Brief} (Goethe-Universität in
Kooperation mit dem Klimareferat) mit personalisiertem Link. Treatment-
Zuweisung erfolgt \textbf{conditional auf Klick} (Average Treatment Effect
on the Treated): jeder Klicker wird randomisiert einem von drei Arme der
Informations-Webseite zugewiesen --- Kontrollgruppe (basale Programminfo),
Stadt-Arm (zusätzlich offizielle Statistik) oder Peer-Arm (zusätzlich
Erfahrungsbericht eines branchen- und größenmatched Peer-Unternehmens).
Outcomes werden auf der Webseite, durch Klick-Tracking sowie über
nachgelagerte Verknüpfung mit den Antrags-Daten der Stadt und dem
Mannheimer Unternehmenspanel erhoben.

% =====================================================================
\section{Stichprobe und Anschreiben}
% =====================================================================

\paragraph{Stichprobenrahmen.} Etwa 20.000 Frankfurter Unternehmen,
gezogen aus dem Mannheimer Unternehmenspanel; ausgeschlossen werden
Unternehmen, die bereits einen Klimabonus-Antrag gestellt haben (nach
Stadt-Admin-Daten).

\paragraph{Anschreiben.} Identischer, treatment-neutraler Brief
(KOMA-Script Vorlage in \var{outreach/letter.tex}) und parallele E-Mail-
Version (\var{outreach/email.tex}). Inhaltlich werden lediglich die
Existenz des Programms und die Eckdaten (bis zu 100.000~EUR pro Maßnahme,
21 Mio Budget, Liste der geförderten Maßnahmen ohne Statistik) genannt.
Studie wird als gemeinsames Forschungsprojekt der Goethe-Universität und
des Klimareferats deklariert; eine kurze Datenschutzklausel weist auf die
Forschungsnutzung hin. Brief und E-Mail enthalten einen \textbf{QR-Code
und einen Kurzlink} (z.\,B. \texttt{klimabonus-frankfurt.de/[ID]}), der
unternehmens-individuell ist.

\paragraph{Erwartete Klickrate.} Konservativ angenommen rund
\textbf{10\,\% Response} $\Rightarrow$ ca.~2.000 Klicker $\Rightarrow$
ca.~660 Beobachtungen pro Treatment-Arm.

\paragraph{Identifikationsstrategie.} ATT, nicht ITT. Die Treatment-
Variation existiert erst nach dem Klick. Wir argumentieren auf der
Population der \emph{marginal-engagierten Unternehmen} --- diejenigen,
die durch Informationsinterventionen erreichbar sind. Ein Stage-1-Modell
\(P(\text{click} \mid \text{MUP-Kovariaten})\) charakterisiert die Klicker
gegen die Frankfurter SME-Population (externe Validität).

% =====================================================================
\section{Treatment-Arme}
% =====================================================================

Alle Arme zeigen die identische Webseite mit Ausnahme einer einzelnen
oberen Sektion (\enquote{Treatment Block}). Alle übrigen Seitenbestandteile
(Hero, Wirtschaftlichkeits-Karten, Förderhöhen-Tabelle, Antragsablauf-
Akkordeon, Kontaktbereich, Anschreiben-Logos) sind \textbf{strikt identisch}
über alle drei Arme.

\subsection*{Arm 1 --- Control (Basisinformation)}
Kurzer neutraler Absatz, der die Komplementarität von Klimaschutz und
wirtschaftlicher Investition skizziert. Keine Statistiken, kein
Erfahrungsbericht.

\subsection*{Arm 2 --- City (T1, offizielle Statistik)}
Stat-Box mit drei prominent dargestellten Zahlen, attribuiert dem
Klimareferat der Stadt Frankfurt:
\begin{itemize}
    \item \textbf{96\,\%} der Anträge werden bewilligt
    \item \textbf{15~Min.} durchschnittlich für den Antrag
    \item \textbf{8.776,97~\euro} durchschnittliche Förderzusage je Antrag
\end{itemize}
Darunter sekundäre Stats: ø~11~Tage Bearbeitungszeit Förderantrag,
ø~32~Tage Bearbeitungszeit Auszahlungsantrag.

\subsection*{Arm 3 --- Peer (T2, peer-bundled testimonial)}
Identische Stat-Box wie in Arm~2 \textbf{plus} darüber eine Peer-Quote-Card:
Foto, Zitat in der Ich-Form, Name, Branche, Größenklasse eines
\emph{vergleichbaren} Frankfurter Unternehmens. Sechs Peer-Profile sind in
\var{klimabonus/peers.py} angelegt; das Matching Empfänger~$\rightarrow$~Peer
erfolgt offline (R/Stata) entlang Branche $\times$ Größenklasse und wird
in der \var{assignment.csv} kodiert (Variable \var{peer\_id}~$\in \{1,\ldots,6\}$).



% =====================================================================
\section{Webseiten-Ablauf}
% =====================================================================

Nach Klick auf den personalisierten Link durchläuft die/der
Teilnehmer:in einen festen, treatment-konstanten Seitenablauf:

\begin{enumerate}[label=\textbf{\arabic*.}]
    \item \textbf{Welcome-Splash} --- automatisches Forwarding,
        Cobranding-Logo, ca.~50~ms sichtbar.
    \item \textbf{Landing-Page} --- enthält den Treatment-Block sowie alle
        identischen Sektionen (Hero, Wirtschaftlichkeits-Karten,
        Förderhöhen-Tabelle, Antragsablauf-Akkordeon, Kontaktbereich,
        Survey-CTA, Direkt-Antrag-CTA). Cookie-Banner erscheint nach 8~Sek;
        Survey-Prompt-Modal nach 2~Min, falls die Person noch keinen
        Survey-CTA geklickt hat.
    \item \textbf{Outcomes-Page} --- Forschungs-Einwilligung (\textbf{required}),
        Bewerbungs-Likelihood-Slider, Checkbox-Grid relevanter
        Fördermaßnahmen, Checkbox-Grid wahrgenommener Hindernisse.
    \item \textbf{Beliefs-Page} --- sieben Items (siehe Variablenliste).
        Alle Beliefs sind \textbf{self-referent} formuliert (\enquote{Ihr
        Antrag / Ihr Unternehmen}). Schließt mit einer Meta-Frage zur
        Confidence in den eigenen Einschätzungen.
    \item \textbf{Commitment-Page} --- Position in der Firma, Wunsch nach
        Folgekontakt (E-Mail, Veranstaltung, Rückruf) und Einwilligung zur
        Datenweitergabe an das Klimareferat (nur sichtbar, wenn
        mindestens ein Kontaktwunsch \enquote{Ja} ist).
    \item \textbf{Abschluss-Page} --- Antragsportal-CTA (extern, in neuem
        Tab), Freitext für Feedback zu städtischen Förderungen allgemein,
        Submit-Button.
    \item \textbf{End-Page} --- identische Darstellung wie Landing
        (mit demselben Treatment-Block), aber ohne Survey-CTA. Erlaubt
        Wiederlesen der Informationen unbegrenzt.
\end{enumerate}

% =====================================================================
\section{Variablenliste}
% =====================================================================

Vollständige Liste aller erhobenen Felder, gegliedert nach Klassen.
Alle Felder werden im \var{custom\_export}-CSV exportiert (oTree Admin
\(\rightarrow\) Export \(\rightarrow\) Custom export). Variablen-Namen
entsprechen den Spaltennamen im Export.




\subsection{Hauptoutcomes (Self-Report, Outcomes-Page)}

\begin{longtable}{@{}p{3.6cm} p{2.0cm} p{9.5cm}@{}}
\toprule
\textbf{Variable} & \textbf{Typ} & \textbf{Beschreibung} \\
\midrule
\endhead
\var{application\_likelihood} & int 0--10  & \enquote{Wie wahrscheinlich ist es, dass Sie in den nächsten 12~Monaten einen Klimabonus-Antrag stellen?} (Slider) \\
\midrule
\multicolumn{3}{@{}l@{}}{\textbf{Relevante Maßnahmen (Mehrfachauswahl, randomisierte Reihenfolge)}} \\
\var{measure\_solar}              & bool & Photovoltaik / Solarthermie \\
\var{measure\_solar\_green\_roof} & bool & Solar-Gründächer \\
\var{measure\_battery}            & bool & Batteriespeicher (mit neuer PV) \\
\var{measure\_greening}           & bool & Dach-, Fassaden- und Hofbegrünung \\
\var{measure\_rainwater}          & bool & Regenwasserspeicher \\
\var{measure\_drinking\_fountain} & bool & Trinkbrunnen \\
\var{measure\_already\_done}      & bool & Ähnliche Maßnahmen bereits umgesetzt (Pioneer-Indikator) \\
\var{measure\_other}              & bool & Sonstiges \\
\var{measure\_other\_text}        & str  & Freitext zu \enquote{Sonstiges} \\
\var{measure\_none}               & bool & Keine davon \\
\var{measures\_order}             & str  & Anzeigereihenfolge (Komma-getrennt) \\
\midrule
\multicolumn{3}{@{}l@{}}{\textbf{Hindernisse (Mehrfachauswahl, randomisierte Reihenfolge)}} \\
\var{barrier\_time}               & bool & Zu hoher Zeitaufwand \\
\var{barrier\_complexity}         & bool & Zu komplizierter Antrag \\
\var{barrier\_uncertainty}        & bool & Unsicherheit, ob bewilligt \\
\var{barrier\_amount}             & bool & Fördersumme zu gering \\
\var{barrier\_property}           & bool & Kein passendes Gebäude / Grundstück \\
\var{barrier\_priority}           & bool & Thema hat keine Priorität \\
\var{barrier\_already\_applied}   & bool & Fördermittel bereits beantragt / genutzt (Screening, kein Hindernis i.\,e.\,S.) \\
\var{barrier\_other}              & bool & Sonstiges \\
\var{barrier\_other\_text}        & str  & Freitext zu \enquote{Sonstiges} \\
\var{barriers\_order}             & str  & Anzeigereihenfolge \\
\bottomrule
\end{longtable}

\subsection{Beliefs (Beliefs-Page, alle self-referent)}

\begin{longtable}{@{}p{3.8cm} p{2.0cm} p{9.3cm}@{}}
\toprule
\textbf{Variable} & \textbf{Typ} & \textbf{Beschreibung} \\
\midrule
\endhead
\var{belief\_approval\_rate}      & int 0--100   & \enquote{Wie wahrscheinlich ist es, dass Ihr Klimabonus-Antrag bewilligt würde?} (Prozent-Slider) \\
\var{belief\_funding\_amount}     & int (EUR)    & \enquote{Wie hoch wäre die Förderung für Ihr Unternehmen?} (Freie Zahleneingabe, 0--100.000) \\
\var{belief\_effort}              & ord. 1--5    & Zeitaufwand für das Ausfüllen Ihres Antrags ($<\!30\,$min bis $>\!1\,$Arbeitstag) \\
\var{belief\_processing\_time}    & ord. 1--5    & Zeit bis Bescheid Ihres Antrags ($<\!2\,$W bis $>\!6\,$M) \\
\var{belief\_payout\_effort}      & ord. 1--5    & Aufwand für Nachweis-Einreichung nach Umsetzung ($<\!30\,$min bis $>\!1\,$AT) \\
\var{belief\_hassle}              & ord. 1--5    & Erwartete Gesamtbelastung des Antragsprozesses (\enquote{Wie belastend\ldots}) \\
\var{belief\_confidence}          & int 1--10    & Meta-Frage am Ende: Sicherheit in den eigenen obigen Einschätzungen (Posterior-Precision-Proxy) \\
\bottomrule
\end{longtable}

\subsection{Befragte:r-Charakteristik, Commitment-Ladder, Feedback}

\begin{longtable}{@{}p{3.6cm} p{2.0cm} p{9.5cm}@{}}
\toprule
\textbf{Variable} & \textbf{Typ} & \textbf{Beschreibung} \\
\midrule
\endhead
\var{respondent\_position}        & cat. & Position in der Firma (Geschäftsführung, Bereichsleitung, kaufm., techn., MA, Sonstige) \\
\var{respondent\_position\_other} & str  & Freitext bei \enquote{Sonstige Position} \\
\var{wants\_email}                & bool & Wunsch nach Info-E-Mail \\
\var{email\_address}              & str  & E-Mail (bedingt erforderlich, wenn \var{wants\_email} oder \var{wants\_event}) \\
\var{wants\_event}                & bool & Wunsch nach Einladung zur Klimabonus-Veranstaltung \\
\var{wants\_callback}             & bool & Wunsch nach Rückruf vom Klimareferat \\
\var{phone\_number}               & str  & Telefonnummer (bedingt erforderlich) \\
\var{feedback\_subsidies}         & str  & Freitext-Feedback zu städtischen Förderungen allgemein \\
\bottomrule
\end{longtable}

\subsection{Revealed Actions (Klick-Tracking)}

\begin{longtable}{@{}p{3.8cm} p{2.0cm} p{9.3cm}@{}}
\toprule
\textbf{Variable} & \textbf{Typ} & \textbf{Beschreibung} \\
\midrule
\endhead
\var{clicked\_portal\_landing}        & bool  & \enquote{Direkt zum Antrag}-Link auf der Landing-Seite geklickt (Fast Path, vor Survey) \\
\var{clicked\_application\_portal}    & bool  & \enquote{Jetzt Klimabonus beantragen}-Button auf Abschluss geklickt (Slow Path, nach Survey) \\
\bottomrule
\end{longtable}

% =====================================================================
\section{Externe Datenquellen (Post-Field)}
% =====================================================================

\subsection{Admin-Daten der Stadt Frankfurt (Klimareferat)}

\begin{itemize}
    \item Tatsächliche Antragsstellung binnen 6 / 12 / 24~Monaten nach Field
    \item Bewilligungsstatus
    \item Förderzusage in Euro (granted amount)
    \item Maßnahmenkategorie (PV, Begrünung, Speicher, \ldots)
    \item Zeitstempel: Antragseinreichung, Bescheid, Auszahlungsantrag, Festsetzung
    \item Tatsächlich ausgezahlte Fördersumme (Auszahlungsantrag)
\end{itemize}

\subsection{Mannheimer Unternehmenspanel (MUP)}

Pre-Treatment-Charakteristik für alle 20.000 angeschriebenen Firmen.
Geplante Kovariaten:

\begin{itemize}
    \item TBD
\end{itemize}



% =====================================================================
\section{Hypothesen-Skizze (vorläufig)}
% =====================================================================

\paragraph{Identifikationsobjekt.} Average Treatment Effect on the Treated
(ATT). Alle Effekte sind interpretiert auf der Subpopulation der
Letter-Klicker (geschätzt ${\sim}\,2.000$).



% =====================================================================



\end{document}
